\section{Models, methods and simulations}
\label{sec:methods}

\para{Physics-based NWP.} An NWP model solves the equations of motion, mass, energy and moisture on a global or regional grid, starting from an analysis built by data assimilation (4D-Var at ECMWF, ensemble Kalman filters in DWD's ICON system). Processes smaller than a grid box, such as radiation, cloud microphysics, convection and boundary-layer turbulence, are parameterised, and for energy these schemes are where most errors originate. Ensembles rerun the model from perturbed states to describe uncertainty. Regional models at 1--3\,km resolve convection explicitly. Table~\ref{tab:models} lists the main systems available as open data.

% Table 5: operational models. Source data: data/literature/nwp_models.csv
\begin{table*}[t]
\centering\small
\caption{\textbf{Operational weather models available as open data} (status November 2026). Grid spacing is approximate; open-data products are sometimes coarser than the native model. Source data: \texttt{data/literature/nwp\_models.csv}.}
\label{tab:models}
\setlength{\tabcolsep}{4pt}
\begin{tabularx}{\textwidth}{@{}L{3.2cm}L{1.9cm}L{2.7cm}L{2.0cm}L{1.6cm}L{1.5cm}Y@{}}
\toprule
Model & Centre & Type & Grid & Range & Runs/day & Licence \\
\midrule
IFS HRES and ENS & ECMWF & physics; ENS 51 members & 9\,km (open 0.25$^\circ$) & 10 / 15 days & 4 & CC-BY-4.0 \\
AIFS Single and ENS & ECMWF & machine learning & about 0.25$^\circ$ & 15 days & 4 & CC-BY-4.0 \\
Extended range & ECMWF & physics ensemble & ensemble & 46 days & 1 & CC-BY-4.0 (partly open) \\
GFS and GEFS & NOAA & physics; GEFS 31 members & 0.25$^\circ$ & 16 days & 4 & public domain \\
AIGFS, AIGEFS, HGEFS & NOAA & machine learning and hybrid & about 0.25$^\circ$ & 16 days & 4 & public domain \\
ICON, ICON-EU & DWD & physics, with ensembles & 13\,km, 6.5\,km & 7.5, 5 days & 4, 8 & open, with attribution \\
ICON-D2 (EPS, RUC) & DWD & convection-permitting & about 2\,km & 2 days (RUC 14\,h) & 8 (RUC 24) & open, with attribution \\
MOSMIX & DWD & statistical post-processing & 5,400 stations & 10 days & several & open, with attribution \\
AROME, ARPEGE & M\'et\'eo-France & physics & 1.3\,km, 5\,km & 2, 4 days & several & Etalab open licence \\
HRRR & NOAA & rapid refresh & 3\,km (USA) & 18--48\,h & 24 & public domain \\
CAMS global & Copernicus & aerosols and dust & about 40\,km & 5 days & 2 & Copernicus licence \\
\bottomrule
\end{tabularx}
\end{table*}


\para{Machine-learning weather models.} Since 2023, models trained on decades of ERA5 reanalysis~\cite{hersbach2020era5} match or beat physics models on many headline scores while running in minutes on a single accelerator: GraphCast~\cite{lam2023graphcast}, GenCast~\cite{price2025gencast}, Aurora~\cite{bodnar2025aurora} and ECMWF's AIFS~\cite{lang2024aifs}, whose deterministic and ensemble versions became operational in 2025~\cite{ecmwf2025aifsens}. NOAA followed in December 2025 with AIGFS, AIGEFS and a 62-member hybrid physics-AI ensemble~\cite{wmo2025noaaai}, and Google's WeatherNext~2 generates large ensembles with hourly output~\cite{google2025weathernext2}. For energy users three caveats matter. Many models output only a subset of variables, so surface radiation, cloud cover and 100\,m wind must be checked; most run at about 0.25$^\circ$ with 6-hourly steps; and deterministic versions smooth extremes at longer lead times. WeatherBench~2 provides common verification~\cite{rasp2024weatherbench2}. The case study (Section~\ref{sec:case}) compares AIFS with IFS and GFS at the station level.

\para{Nowcasting.} For the first hours, observations beat models. Geostationary satellites (Meteosat Third Generation has scanned Europe operationally since December 2024) give cloud fields every few minutes, which are converted to irradiance and advected with cloud motion vectors or learned models; sky cameras cover the next minutes at a single plant; radar tracks precipitation and convective cells; and rapid-update NWP such as ICON-D2-RUC refreshes hourly. Operational systems blend these sources with live plant data, shifting weight from measurements to NWP as the lead time grows.

\para{From weather to power.} The physical chain of Section~\ref{sec:foundations} is implemented in pvlib~\cite{holmgren2018pvlib} for PV and in windpowerlib~\cite{windpowerlib}, PyWake or FLORIS for wind. The statistical layer ranges from linear MOS through gradient-boosted trees, the workhorse for tabular NWP features, to neural networks on NWP fields and graph networks across many sites. Typical inputs are several NWP sources, clear-sky irradiance and solar angles, lead time, calendar terms, ensemble statistics and, for intraday use, recent measurements; targets are normalised by capacity or clear-sky power. Portfolio and regional forecasts upscale from reference plants and reconcile plant, portfolio and regional forecasts so that they add up.

\para{Probabilistic forecasts.} Raw ensembles passed through a power model are physically consistent but usually under-dispersive, so they are calibrated with EMOS or Bayesian model averaging~\cite{gneiting2005emos}. Quantile regression and quantile forests learn quantiles directly; distributional models predict full densities; analog ensembles reuse what happened after similar past forecasts; conformal prediction adds calibrated intervals to any model. Storage optimisation and stochastic dispatch need scenarios that respect temporal correlation, typically generated with a Gaussian copula over marginal quantiles.

\para{Combining models.} Because NWP systems fail differently, blending several sources usually beats the best single source, with static, lead-dependent or performance-based weights. Proposition~\ref{prop:combine} states when even an equal-weight average of two models helps.

\para{Simulations.} Beyond forecasting, the industry relies on a family of simulations that use the same weather data: energy-yield simulation (PVsyst, SAM, pvlib, WAsP, WindPRO); wake and flow simulation from engineering models to computational fluid dynamics and large-eddy simulation; mesoscale downscaling with WRF or ICON for wind atlases; Monte Carlo propagation of uncertainty to P90 and revenue at risk; resampling of 30 or more historical weather years for budgets, PPAs and adequacy studies; weather-window simulation for construction and offshore logistics; hydrological rainfall-runoff and snowmelt models; icing, snow and soiling loss models; power-system and market models such as PyPSA or Antares for price scenarios; and linear or mixed-integer optimisation of storage dispatch under forecast scenarios.
